\chapter{API RESTful proposta: requisitos, análise, design e implementação}

\section{Visão geral da solução}

Este capítulo apresenta a solução tecnológica desenvolvida para apoiar o processo de captação e gerenciamento de solicitações de doação de sangue. A proposta materializa-se na construção de uma API RESTful focada exclusivamente na camada de \textit{backend}, responsável pelo gerenciamento de participantes (pessoas e organizações), solicitações de doação, registro do histórico de doadores e recomendação de solicitações compatíveis.

O desenvolvimento foi orientado pelos princípios do DDD e pelos princípios SOLID, já fundamentados no Capítulo~2. Essa abordagem permitiu que a estrutura do software refletisse a realidade do domínio da hemoterapia, reduzindo o acoplamento entre os componentes e facilitando a evolução contínua do sistema. As seções seguintes detalham o levantamento de requisitos, a modelagem do sistema, a organização do domínio e da arquitetura, a implementação técnica e, por fim, a estratégia de recomendação e de cumprimento de metas, que constitui o núcleo funcional da API.

\section{Levantamento de requisitos}

A extração de requisitos ocorreu a partir da análise documental e bibliográfica descrita no Capítulo~4, alinhando as necessidades dos diferentes atores envolvidos no processo de doação de sangue às capacidades técnicas do sistema proposto.

Os requisitos funcionais representam as operações que a API disponibiliza para atender às rotinas de cadastro, gerenciamento de solicitações, recomendação e controle de doações. A tabela~\ref{tab:rf} sintetiza as funcionalidades identificadas e implementadas, refletindo diretamente os recursos (\textit{endpoints}) expostos pela API.

\begin{table}[H]
  \centering
  \caption{Requisitos funcionais}
  \label{tab:rf}
  \small
  \begin{tabularx}{\textwidth}{c X}
    \toprule
    \textbf{Código} & \textbf{Descrição} \\
    \midrule
    RF01 & Permitir o cadastro de pessoas no sistema. \\
    RF02 & Permitir o cadastro de organizações no sistema. \\
    RF03 & Permitir que uma pessoa cadastrada registre um perfil de doador. \\
    RF04 & Permitir que uma pessoa ou organização registre um perfil de requisitante. \\
    RF05 & Permitir que uma organização cadastrada registre um perfil de banco de sangue (hemocentro). \\
    RF06 & Permitir a autenticação de usuários por meio de credenciais válidas, emitindo um JWT. \\
    RF07 & Permitir a atualização dos dados cadastrais de um participante (\textit{party}). \\
    RF08 & Permitir a atualização do perfil e das preferências de recomendação do doador. \\
    RF09 & Permitir que requisitantes criem solicitações de doação de sangue. \\
    RF10 & Permitir a atualização dos dados de uma solicitação (meta de bolsas e/ou data-limite). \\
    RF11 & Permitir o cancelamento de uma solicitação de doação. \\
    RF12 & Permitir a consulta das solicitações de doação de um participante. \\
    RF13 & Permitir a consulta das solicitações de doação recomendadas para um doador. \\
    RF14 & Permitir a notificação de doadores potenciais elegíveis para uma solicitação. \\
    RF15 & Permitir o registro de uma doação, informando a data de intenção (pendente) ou a data efetiva (concluída). \\
    RF16 & Permitir o reagendamento de uma doação pendente. \\
    RF17 & Permitir a confirmação (conclusão) de uma doação pendente. \\
    RF18 & Permitir a consulta do histórico de doações de um doador. \\
    RF19 & Permitir a consulta do resumo (\textit{summary}) de um doador. \\
    \bottomrule
  \end{tabularx}
  \fonte{Elaborado pelos próprios autores (2026).}
\end{table}


Ressalta-se que a lista de requisitos funcionais apresentada não tem a pretensão de esgotar todas as rotinas administrativas de um hemocentro ou unidade hospitalar, mas de delimitar um conjunto mínimo e coeso de operações indispensáveis para validar os conceitos teóricos e arquiteturais propostos. Esse escopo abrange o ciclo fundamental da doação de sangue ao contemplar, de maneira integrada, a gestão da demanda (abertura e acompanhamento de solicitações por unidades de saúde e pacientes), a oferta (disponibilidade, elegibilidade e intenções de doação pelos voluntários) e o registro histórico das doações realizadas.


Além das funcionalidades diretas, foram definidos requisitos não funcionais para estabelecer os parâmetros de qualidade, segurança e arquitetura esperados. Considerando a adoção do DDD e da arquitetura em camadas, a Tabela~\ref{tab:rnf} apresenta as diretrizes técnicas estabelecidas para a solução.

\begin{table}[H]
  \centering
  \caption{Requisitos não funcionais}
  \label{tab:rnf}
  \small
  \begin{tabularx}{\textwidth}{c X}
    \toprule
    \textbf{Código} & \textbf{Descrição} \\
    \midrule
    RNF01 & O sistema deve disponibilizar seus serviços por meio de uma API RESTful utilizando o protocolo HTTP. \\
    RNF02 & O sistema deve autenticar usuários por meio de JWT. \\
    RNF03 & O sistema deve persistir os dados em banco de dados NoSQL (MongoDB). \\
    RNF04 & O sistema deve permitir evolução e manutenção sem impacto significativo nas demais camadas da aplicação. \\
    RNF05 & O sistema deve seguir uma arquitetura em camadas orientada pelos princípios do DDD. \\
    RNF06 & O sistema deve aplicar os princípios SOLID para favorecer baixo acoplamento e alta coesão. \\
    RNF07 & A API deve retornar respostas em formato JSON. \\
    RNF08 & O sistema deve validar os dados recebidos antes do processamento das operações. \\
    RNF09 & O sistema deve controlar o acesso às funcionalidades protegidas por meio de autorização baseada em perfis (\textit{roles}) de usuário. \\
    RNF10 & O sistema deve preservar a consistência de escritas concorrentes sobre uma mesma solicitação de doação, evitando a perda de atualizações. \\
    \bottomrule
  \end{tabularx}
  \fonte{Elaborado pelos próprios autores (2026).}
\end{table}


Diferentemente dos requisitos funcionais, que especificam os serviços e as ações que o sistema deve executar, os requisitos não funcionais definem as qualidades, restrições e propriedades estruturais do \textit{software}. Nesse sentido, os itens elencados estabelecem parâmetros de segurança (como autenticação sem estado e autorização por papéis), manutenibilidade e baixo acoplamento (assegurados pela divisão em camadas orientada ao DDD e ao SOLID), integridade de dados (por meio do controle de concorrência otimista) e padronização das comunicações (via API RESTful e formato JSON), sustentando a robustez técnica da solução desenvolvida.


\section{Modelagem do sistema}

Após a consolidação dos requisitos, procedeu-se com a modelagem do sistema para descrever as interações entre os atores e a API, permitindo visualizar o comportamento esperado da solução sob a perspectiva estrutural e de uso.

\subsection{Casos de uso}

A modelagem de casos de uso, em notação UML, sintetiza as operações disponíveis para os diferentes perfis de acesso. A Tabela~\ref{tab:uc} apresenta o mapeamento dos casos de uso e seus respectivos atores, incluindo tanto as operações representadas nos diagramas a seguir quanto operações complementares oferecidas pela API para os mesmos atores.

\begin{table}[H]
  \centering
  \caption{Lista de casos de uso}
  \label{tab:uc}
  \small
  \begin{tabularx}{\textwidth}{c >{\raggedright\arraybackslash}X >{\raggedright\arraybackslash}p{3.2cm}}
    \toprule
    \textbf{Código} & \textbf{Caso de uso} & \textbf{Ator principal} \\
    \midrule
    UC01 & Cadastrar-se como pessoa (\textit{Sign Up as Person}) & Guest \\
    UC02 & Cadastrar-se como organização (\textit{Sign Up as Organization}) & Guest \\
    UC03 & Autenticar (\textit{Authenticate}) & Guest \\
    UC04 & Registrar perfil de doador & Person \\
    UC05 & Registrar perfil de requisitante & Person / Organization \\
    UC06 & Registrar perfil de banco de sangue & Organization \\
    UC07 & Atualizar dados cadastrais (\textit{Update Party}) & Person / Organization \\
    UC08 & Atualizar perfil e preferências do doador (\textit{Update Donor}) & Donor \\
    UC09 & Registrar solicitação de doação (\textit{Register Donation Request}) & Requester \\
    UC10 & Consultar recomendações de solicitações (\textit{Consult Recommendations}) & Donor \\
    UC11 & Registrar doação (\textit{Register Donation}) & Donor \\
    UC12 & Completar doação pendente (\textit{Complete a Pending Donation}) & Donor \\
    UC13 & Consultar histórico de doações (\textit{See Donation History}) & Donor \\
    UC14 & Gerenciar solicitações --- atualizar dados, cancelar e consultar & Requester \\
    UC15 & Notificar doadores potenciais elegíveis & Requester \\
    UC16 & Reagendar doação pendente & Donor \\
    UC17 & Consultar resumo do doador (\textit{Donor Summary}) & Donor \\
    \bottomrule
  \end{tabularx}
  \fonte{Elaborado pelos próprios autores (2026).}
\end{table}

A representação gráfica dessas interações é dividida de acordo com o nível de acesso dos atores. A Figura~\ref{fig:uc-access-baseProfiles} apresenta o diagrama de caso de uso de acesso e perfis base (\textit{Guest, Person and Organization}).

\begin{figure}[H]
  \centering
  \caption{Diagrama de caso de uso de acesso e perfis base}
  \includegraphics[width=0.6\textwidth]{figura03-uc-acesso-perfisBase}
  \fonte{Elaborado pelos próprios autores (2026).}
  \label{fig:uc-access-baseProfiles}
\end{figure}

O ator \textit{Guest} representa utilizadores sem cadastro prévio, limitados às operações de criação de contas e autenticação (UC01, UC02 e UC03). Os atores \textit{Person} e \textit{Organization} representam os usuários recém-cadastrados, permitindo a atualização de dados da conta (UC07) e a definição de seus papéis específicos no sistema de domínio, como Doador (UC04), Requisitante (UC05) ou Banco de Sangue (UC06).

Já a Figura~\ref{fig:uc-requester} representa o diagrama de caso de uso do ator (\textit{Requester}).

\begin{figure}[H]
  \centering
  \caption{Diagrama de caso de uso de requisitante}
  \includegraphics[width=0.65\textwidth]{figura04-uc-requisitante}
  \fonte{Elaborado pelos próprios autores (2026).}
  \label{fig:uc-requester}
\end{figure}

O ator \textit{Requester} concentra as permissões relacionadas à criação de solicitações de doação (UC09) e à sua gestão (UC14), além do disparo de notificações a doadores potenciais quando uma solicitação necessita de maior alcance (UC15).

Por fim, a Figura~\ref{fig:uc-donor} representa o diagrama de caso de uso do ator (\textit{Donor}).

\begin{figure}[H]
  \centering
  \caption{Diagrama de caso de uso de doador}
  \includegraphics[width=0.7\textwidth]{figura05-uc-doador}
  \fonte{Elaborado pelos próprios autores (2026).}
  \label{fig:uc-donor}
\end{figure}

O ator \textit{Donor} possui acesso ao núcleo do processo de doação, incluindo a atualização de seu perfil clínico (UC08), a consulta às solicitações recomendadas (UC10), o registro unificado de doações (UC11), a conclusão (UC12) e o reagendamento de doações pendentes (UC16), a consulta ao próprio histórico hemoterápico (UC13) e a emissão de seu resumo (UC17). Vale ressaltar que, no UC11, a doação nasce pendente quando o doador informa a data de intenção ou concluída quando informa a data efetiva --- um único caso de uso, sem fábricas concorrentes no domínio.

\subsubsubsection{Detalhamento dos casos de uso}

Devido à abrangência da API e visando manter a concisão documental, optou-se por detalhar apenas os casos de uso críticos para o núcleo da aplicação, adotando a estrutura formal de especificação textual proposta por Cockburn (2001). Essa abordagem documenta o fluxo de eventos, as restrições operacionais e as pré-condições de segurança exigidas, justificando arquiteturalmente a ausência de relacionamentos visuais de dependência de autenticação nos diagramas UML. A Tabela~\ref{tab:detalhamento-uc} sintetiza essas especificações.

\begin{table}[H]
  \centering
  \caption{Especificação dos casos de uso críticos}
  \label{tab:detalhamento-uc}
  \small
  \begin{tabularx}{\textwidth}{>{\bfseries}p{3.3cm} X}
    \toprule
    \multicolumn{2}{l}{\textbf{UC09: Registrar solicitação de doação}} \\
    \midrule
    Ator principal & Requisitante (\textit{Requester}) \\
    Pré-condição & Usuário autenticado com perfil de Requisitante ativo. \\
    Fluxo principal & 
    1. O Requisitante solicita a criação de uma nova campanha.\\
    & 2. Informa dados da necessidade (tipo sanguíneo, bolsas, prazo, urgência e hemocentro).\\
    & 3. O sistema valida as regras de negócio e os invariantes de domínio.\\
    & 4. O sistema persiste a solicitação no estado ativo.\\
    & 5. O sistema retorna a confirmação com o identificador único gerado. \\
    Cenário alternativo & \textbf{3a.} Falha de validação (prazo retroativo ou meta nula): rejeita a operação e retorna mensagem de erro estruturada. \\
    \midrule
    \multicolumn{2}{l}{\textbf{UC10: Consultar recomendações de solicitações}} \\
    \midrule
    Ator principal & Doador (\textit{Donor}) \\
    Pré-condição & Usuário autenticado com perfil de Doador ativo. \\
    Fluxo principal & 
    1. O Doador solicita a consulta de campanhas recomendadas.\\
    & 2. O sistema avalia a elegibilidade temporal com base na última doação.\\
    & 3. O sistema recupera solicitações ativas e avalia a compatibilidade tipológica.\\
    & 4. O sistema aplica o filtro geoespacial pela distância máxima configurada.\\
    & 5. O sistema ordena por urgência, proximidade e prazo, retornando a listagem. \\
    Cenário alternativo & \textbf{2a.} Doador temporariamente inelegível: encerra o processamento e retorna listagem vazia. \\
    \midrule
    \multicolumn{2}{l}{\textbf{UC12: Completar doação pendente}} \\
    \midrule
    Ator principal & Doador (\textit{Donor}) \\
    Pré-condição & Usuário autenticado com perfil de Doador ativo. \\
    Fluxo principal & 
    1. O Doador submete a data de efetivação de intenção previamente registrada.\\
    & 2. O sistema recupera a doação pendente associada ao identificador.\\
    & 3. O sistema valida a consistência cronológica em relação à data corrente.\\
    & 4. O sistema transiciona o estado do agregado para concluído.\\
    & 5. O sistema atualiza o registro do doador, recalculando a janela de descanso regulamentar. \\
    Cenário alternativo & \textbf{3a.} Data inconsistente (data futura): rejeita a atualização e mantém a doação pendente. \\
    \bottomrule
  \end{tabularx}
  \fonte{Elaborado pelos próprios autores (2026).}
\end{table}

\subsection{Diagramas de classes}

Para representar a estrutura estática da aplicação e evidenciar a aplicação do DDD na prática --- cujos conceitos de entidade, agregado e objeto de valor já foram apresentados no Capítulo~2 ---, a modelagem de classes foi concebida em duas perspectivas distintas: conceitual e de implementação. A perspectiva conceitual mapeia o vocabulário do negócio e as relações semânticas do mundo real, omitindo métodos e tipos de dados para facilitar a validação com especialistas. Em contrapartida, a perspectiva de implementação atua como uma representação estrutural do código-fonte construído, detalhando os limites transacionais, o encapsulamento e as diretrizes arquiteturais.

A Figura~\ref{fig:classe-participantes} apresenta o diagrama de classes conceitual focado nos participantes e em seus papéis (\textit{Roles}).

\begin{figure}[H]
  \centering
  \caption{Diagrama conceitual dos participantes}
  \includegraphics[width=0.65\textwidth]{figura06-classes-participantes}
  \fonte{Elaborado pelos próprios autores (2026).}
  \label{fig:classe-participantes}
\end{figure}

Observa-se a utilização da abstração central \texttt{Party}, que atua como classe base para \texttt{Person} (Pessoa) e \texttt{Organization} (Organização). Essa modelagem reflete a flexibilidade do domínio, permitindo que as entidades assumam múltiplos papéis, como \texttt{Donor} (Doador) e \texttt{Requester} (Requisitante), eliminando a duplicidade de dados no sistema.

Dando prosseguimento à representação do negócio, a Figura~\ref{fig:classe-doacao} ilustra o diagrama de classes conceitual do processo de doação.

\begin{figure}[H]
  \centering
  \caption{Diagrama conceitual do processo de doação}
  \includegraphics[width=0.65\textwidth]{figura07-classes-doacao}
  \fonte{Elaborado pelos próprios autores (2026).}
  \label{fig:classe-doacao}
\end{figure}

Neste nível de abstração, evidencia-se a associação semântica entre a solicitação (\codigo{DonationRequest}) e o ato de doar (\codigo{Donation}). O fluxo natural demonstra que um \codigo{Requester} origina a solicitação, que é vinculada a um \codigo{BloodCenter} e, posteriormente, recebe as doações efetivadas por um \codigo{Donor}.

Elevando o rigor para o design de software, a Figura~\ref{fig:impl-request} detalha a implementação técnica do agregado \texttt{DonationRequest}.

\begin{figure}[H]
  \centering
  \caption{Diagrama de implementação: agregado DonationRequest}
  \includegraphics[width=0.9\textwidth]{figura08-impl-donation-request}
  \fonte{Elaborado pelos próprios autores (2026).}
  \label{fig:impl-request}
\end{figure}

A classe é modelada como uma raiz de agregado (\texttt{<<AggregateRoot>>}), atuando como a única porta de entrada para a proteção de invariantes, como as datas-limite e a urgência da solicitação. Observa-se a aplicação da herança por meio da superclasse abstrata \texttt{DomainObject}, responsável por compor os identificadores únicos globais (\texttt{<<ValueObject>> DomainID}). Destaca-se também a dependência tracejada (\texttt{<<use>>}) com a classe \texttt{Donation}, a qual evidencia o baixo acoplamento da arquitetura: a requisição não armazena doações em memória, apenas as avalia de forma transitória por meio de parâmetros no método \texttt{acceptsDonation()}.

Por fim, a Figura~\ref{fig:impl-donation} apresenta o diagrama de implementação do agregado \texttt{Donation}.

\begin{figure}[H]
  \centering
  \caption{Diagrama de implementação: agregado Donation}
  \includegraphics[width=0.75\textwidth]{figura09-impl-donation}
  \fonte{Elaborado pelos próprios autores (2026).}
  \label{fig:impl-donation}
\end{figure}

Tratada também como uma raiz de agregado autônoma, a classe concentra uma única fábrica \codigo{create} e o ciclo de vida pendente--concluída--cancelada. O estado não é armazenado como flags independentes: deriva das datas de intenção (\codigo{intendedDate}), de efetivação (\codigo{donationDate}) e de cancelamento (\codigo{cancelledAt}). Completar uma doação pendente preenche a data efetiva sem apagar a intenção; reagendar e cancelar permanecem transições sobre o mesmo agregado. A navegabilidade unidirecional explícita demonstra que a doação possui a referência imediata do \texttt{Donor} e do \texttt{BloodCenter} envolvidos, assegurando a consistência da operação sem a necessidade de carregar listas e metadados completos de uma solicitação simultaneamente.

\section{Domínio e arquitetura}

A partir das interações mapeadas, estruturou-se o núcleo da aplicação: o modelo de domínio e a organização das camadas que o protegem de dependências tecnológicas externas.

\subsection{Linguagem ubíqua e organização em camadas}

Conforme estabelecido nos fundamentos do DDD (Capítulo~2), a modelagem foi construída sobre uma linguagem ubíqua, garantindo que termos como \texttt{DonationRequest}, \texttt{Donor}, \texttt{Requester} e \texttt{BloodCenter} transitem com a mesma semântica tanto nas discussões de projeto quanto no código-fonte final. Essa padronização reduz ambiguidades entre a gestão de perfis de participantes e o ciclo de vida das doações. 

A organização estrutural do código adota o padrão orientado a entidades de domínio (\textit{Package-by-Entity} / \textit{Package-by-Feature}), promovendo alta coesão e baixo acoplamento entre os conceitos do sistema. Em vez de uma divisão horizontal global agrupando classes exclusivamente por camada técnica, o sistema é particionado em pacotes de alto nível que representam diretamente os agregados e entidades de domínio (\texttt{request}, \texttt{donation}, \texttt{party}, \texttt{role} e \texttt{auth}), complementados por um pacote compartilhado (\texttt{shared}) para conceitos transversais (\textit{Shared Kernel}).

Internamente, cada pacote de entidade é estruturado nos subpacotes que materializam as camadas arquiteturais clássicas, garantindo a separação estrita de responsabilidades:

\begin{itemize}[leftmargin=*]
  \item \texttt{domain}: entidades, raízes de agregado, objetos de valor, regras de negócio e interfaces de portas de repositórios e serviços de domínio, totalmente livres de acoplamento com frameworks ou tecnologias externas;
  \item \texttt{application}: orquestração dos casos de uso, coordenando o fluxo de dados entre as portas de repositório e os modelos de domínio;
  \item \texttt{infrastructure}: implementações concretas dos repositórios (Spring Data MongoDB), esquemas de persistência de banco de dados, filtros de segurança e clientes de serviços externos;
  \item \texttt{presentation}: controladores REST (\textit{controllers}) e objetos de transferência de dados (DTOs) responsáveis pela exposição HTTP da API, recepção de requisições e validação sintática primária.
\end{itemize}

O fluxo de dados segue os princípios da arquitetura limpa: a camada de apresentação invoca a camada de aplicação, que opera sobre o domínio (\texttt{presentation} $\rightarrow$ \texttt{application} $\rightarrow$ \texttt{domain}). A camada de infraestrutura, por sua vez, implementa as interfaces definidas no domínio (\texttt{infrastructure} $\rightarrow$ \texttt{domain}), garantindo a aplicação do Princípio da Inversão de Dependência (DIP) entre os módulos de alto e baixo nível \cite{martin2008}.

\subsection{Entidades, agregados e objetos de valor}

A estrutura central do domínio foi modelada de forma hierárquica e orientada a objetos. A Tabela~\ref{tab:entidades} sintetiza as principais entidades e agregados responsáveis por garantir a consistência das transações no sistema.

\begin{table}[H]
  \centering
  \caption{Entidades e suas responsabilidades}
  \label{tab:entidades}
  \begin{tabularx}{\textwidth}{>{\raggedright\arraybackslash}p{4cm} X}
    \toprule
    \textbf{Entidade} & \textbf{Responsabilidade} \\
    \midrule
    Party & Abstrai características comuns dos participantes do sistema. \\
    Person & Representa pessoas físicas cadastradas. \\
    Organization & Representa instituições cadastradas. \\
    Donor & Representa o perfil de doador e suas regras de elegibilidade. \\
    Requester & Representa o perfil requisitante. \\
    BloodCenter & Representa o perfil de banco de sangue vinculado a uma organização. \\
    DonationRequest & Representa uma solicitação de doação de sangue, incluindo sua meta e seu progresso. \\
    Donation & Representa uma doação registrada no sistema, cujo estado (pendente, concluída ou cancelada) deriva das datas de intenção, de efetivação e de cancelamento. \\
    \bottomrule
  \end{tabularx}
  \fonte{Elaborado pelos próprios autores (2026).}
\end{table}

A entidade \texttt{Party} foi adotada como abstração de mais alto nível da hierarquia, aplicando herança para derivar \texttt{Person} e \texttt{Organization}. A partir dessas bases, o sistema utiliza polimorfismo para permitir que os participantes assumam papéis (\textit{Roles}) específicos, como \texttt{Donor}, \texttt{Requester} ou \texttt{BloodCenter}, dependendo de suas ações no sistema.

As entidades \texttt{DonationRequest} e \texttt{Donation} operam como raízes de agregados (\textit{Aggregate Roots}), concentrando regras de negócio --- como a validação de datas-limite e a verificação de compatibilidade sanguínea --- e protegendo o estado interno dos objetos por meio do encapsulamento. Adicionalmente, conceitos imutáveis sem identidade própria, como \texttt{Address} (Endereço), \texttt{BloodType} (Tipo Sanguíneo) e \texttt{DomainID}, foram modelados como objetos de valor, na linha proposta por \citeonline{evans2003} e detalhada por \citeonline{vernon2013}, aumentando a segurança e a previsibilidade do modelo.

\section{Implementação e Fluxo de Dados}

A implementação técnica da API orienta-se pelos princípios da arquitetura em camadas e do \textit{Domain-Driven Design} (DDD). Para fins de exemplificar a aplicação prática desses conceitos, esta seção demonstra o fluxo completo de uma entidade --- o registro de uma \textbf{Solicitação de Doação} (\codigo{DonationRequest}) ---, acompanhando a requisição desde a sua chegada no controlador REST até a sua persistência no banco de dados com controle de concorrência. Esse rastreamento horizontal ilustra o papel de cada camada arquitetural e as transformações sofridas pelos dados.

\subsection{Camada de Apresentação: REST Controller e DTOs}

A camada de apresentação atua como o ponto de entrada da aplicação, isolando o núcleo do sistema dos detalhes do protocolo HTTP. O \codigo{CreateDonationRequestController} é responsável por interceptar as requisições, executar validações sintáticas básicas e converter o \textit{payload} da requisição (DTO) para o modelo de entrada do caso de uso.

\begin{lstlisting}[language=Java, caption={Exemplo de REST Controller da Solicitação de Doação}, label={lst:create-request-controller}]
@RestController
@RequestMapping("/donation-requests")
public class CreateDonationRequestController {

  private final CreateDonationRequestUseCase useCase;

  public CreateDonationRequestController(CreateDonationRequestUseCase useCase) {
    this.useCase = useCase;
  }

  @PostMapping
  public ResponseEntity<CreateDonationRequestResponseDto> create(
      @RequestBody CreateDonationRequestDto payload) {
    
    // Validacoes sintaticas primarias
    requireNonNull(payload.goalBloodBags(), "goalBloodBags cannot be null");
    requireSamePartyOrAdmin(payload.partyId()); // Autorizacao

    // Conversao do DTO para o objeto de entrada (Input) do Caso de Uso
    var output = useCase.execute(new Input(
        payload.partyId(),
        payload.organizationId(),
        payload.bloodTypeNeeded(),
        payload.goalBloodBags(),
        payload.dateLimit(),
        payload.urgency(),
        payload.directedTo()));

    return ResponseEntity
        .status(HttpStatus.CREATED)
        .body(CreateDonationRequestResponseDto.from(output));
  }
}
\end{lstlisting}

\subsection{Camada de Aplicação: Casos de Uso}

A camada de aplicação orquestra o fluxo de dados coordenando chamadas aos repositórios e aos serviços de domínio. O \codigo{CreateDonationRequestUseCase} carrega as entidades associadas (\codigo{Requester} e \codigo{BloodCenter}), invoca a fábrica da entidade e instrui o repositório a persistir o novo agregado. Esta camada desconhece qualquer dependência de infraestrutura, utilizando apenas portas de interfaces.

\begin{lstlisting}[language=Java, caption={Exemplo de Caso de Uso de Criação}, label={lst:create-request-usecase}]
@Service
public class CreateDonationRequestUseCase {

  private final DonationRequestRepositoryInterface repository;
  private final RequesterRepositoryInterface requesterRepo;
  private final BloodCenterRepositoryInterface bloodCenterRepo;

  // Construtor omitido para brevidade...

  public Output execute(Input input, LocalDate currentDate) {
    // 1. Parse de Identificadores e Value Objects
    DomainID partyId = DomainIdParser.parse(input.partyId(), "partyId");
    BloodType bloodType = BloodType.of(input.bloodTypeNeeded());

    // 2. Busca das entidades relacionadas nas portas de repositorio
    Requester requester = requesterRepo.findByPartyId(partyId)
        .orElseThrow(() -> new NotFoundException("Requester not found"));
    BloodCenter bloodCenter = bloodCenterRepo.findByPartyId(organizationId)
        .orElseThrow(() -> new NotFoundException("Blood center not found"));

    // 3. Invocacao da fabrica do Agregado de Dominio
    DonationRequest request = DonationRequest.create(
        requester,
        bloodCenter,
        bloodType,
        input.goalBloodBags(),
        input.dateLimit(),
        currentDate,
        parseUrgency(input.urgency()),
        input.directedTo());

    // 4. Persistencia e retorno
    repository.save(request);
    return Output.from(request);
  }
}
\end{lstlisting}

\subsection{Camada de Domínio: Agregado e Regras de Negócio}

O coração do software reside na camada de domínio. A classe \codigo{DonationRequest} atua como a raiz do agregado, protegendo os seus invariantes e garantindo que nenhuma solicitação seja instanciada em um estado inválido. A lógica de negócio essencial (como não permitir metas zeradas ou datas no passado) fica concentrada no método fábrica \codigo{create(...)}.

\begin{lstlisting}[language=Java, caption={Fábrica do Agregado de Domínio}, label={lst:donation-request-domain}]
public class DonationRequest extends DomainObject {
  private Requester requester;
  private BloodCenter bloodCenter;
  private BloodType bloodTypeNeeded;
  private int goalBloodBags;
  private LocalDate dateLimit;

  public static DonationRequest create(
      Requester requester, BloodCenter bloodCenter, BloodType bloodTypeNeeded,
      int goalBloodBags, LocalDate dateLimit, LocalDate currentDate,
      Urgency urgency, String directedTo) {
      
    // Protecao de Invariantes Essenciais
    if (goalBloodBags <= 0)
      throw new IllegalArgumentException("Goal blood bags must be > zero");
    if (dateLimit.isBefore(currentDate))
      throw new IllegalArgumentException("Limit date invalid");
      
    return new DonationRequest(
        requester, bloodCenter, bloodTypeNeeded, goalBloodBags,
        dateLimit, currentDate, urgency, directedTo);
  }
}
\end{lstlisting}

\subsection{Camada de Infraestrutura: Repositório e Controle de Concorrência}

Na extremidade da arquitetura, a camada de persistência traduz o modelo conceitual rico do domínio para as estruturas do banco de dados (MongoDB). Os adaptadores implementam as interfaces requeridas pelos casos de uso, mapeando o agregado de domínio para um \textit{Schema} correspondente.

Ademais, para garantir a integridade das transações contra atualizações simultâneas (como alterar simultaneamente o prazo e a meta de uma mesma solicitação), o esquema incorpora um controle de versão otimista suportado pela anotação \codigo{@Version}.

\begin{lstlisting}[language=Java, caption={Mapeamento e Controle de Concorrência Otimista}, label={lst:donation-request-schema}]
@Document(collection = "donation_requests")
public class DonationRequestSchema {

  @Id
  private String id;
  private String requesterId;
  private String bloodTypeNeeded;
  private int goalBloodBags;

  @Version
  private Long version; // Controle de Concorrencia Otimista

  // Metodo de conversao (Dominio -> Infraestrutura)
  public static DonationRequestSchema fromDomain(DonationRequest domain) {
    DonationRequestSchema schema = new DonationRequestSchema();
    schema.id = domain.getId().getValue().toString();
    schema.requesterId = domain.getRequester().getId().getValue().toString();
    schema.goalBloodBags = domain.getGoalBloodBags();
    schema.version = domain.getVersion();
    return schema;
  }
}
\end{lstlisting}

Para preservar a arquitetura, todas as exceções intrínsecas ao banco de dados são tratadas e convertidas por um utilitário interno (\codigo{OptimisticConcurrency}) dentro do próprio adaptador, impedindo o vazamento de exceções tecnológicas de persistência para as camadas superiores:

\begin{lstlisting}[language=Java, caption={Adaptador de Repositório}, label={lst:donation-request-repo}]
@Repository
public class DonationRequestRepositoryImpl implements DonationRequestRepositoryInterface {

  private final DonationRequestMongoRepository mongoRepository;

  @Override
  public void save(DonationRequest request) {
    OptimisticConcurrency.execute(() -> {
      DonationRequestSchema schema = DonationRequestSchema.fromDomain(request);
      mongoRepository.save(schema);
    });
  }
}
\end{lstlisting}

\subsection{Segurança e autenticação}

O controle de acesso é implementado por meio de autenticação baseada em JWT. O endpoint público de login (\texttt{/auth/login}) valida as credenciais do usuário e emite um \textit{token} assinado, que passa a ser exigido nas demais chamadas autenticadas por meio de um filtro (\texttt{JwtAuthenticationFilter}) na camada \texttt{infra/security}. Esse filtro intercepta as requisições antes de chegarem aos \textit{controllers}, valida o \textit{token} e popula o contexto de segurança com a identidade do participante autenticado, permitindo que os casos de uso apliquem regras de autorização --- como a verificação de que apenas o próprio requisitante pode gerenciar ou notificar doadores sobre uma solicitação de sua autoria.

\subsection{Serviços externos e documentação da API}

A notificação de doadores é abstraída por uma interface de domínio (\codigo{NotificationServiceInterface}), permitindo múltiplas implementações na camada de infraestrutura --- como uma implementação por e-mail e uma implementação simplificada para ambiente de desenvolvimento --- sem que o caso de uso que dispara a notificação precise conhecer o mecanismo concreto de envio. Essa inversão de dependência, alinhada aos princípios SOLID discutidos no Capítulo~2, evita que o domínio fique acoplado a bibliotecas de terceiros. 

Para a documentação e a exploração interativa dos recursos da API, foi integrada a biblioteca \textit{springdoc-openapi}, que gera automaticamente a especificação OpenAPI a partir dos \textit{controllers} e disponibiliza uma interface de teste dos \textit{endpoints}. Essa documentação viva reduz o custo de manutenção da documentação técnica e facilita a validação manual dos fluxos descritos neste capítulo.

\section{Recomendação de solicitações e cumprimento de metas}
\label{sec:recomendacao-fulfillment}

Esta seção detalha o mecanismo central da API: a recomendação de solicitações compatíveis a doadores, a notificação de doadores potenciais e o acompanhamento do cumprimento de metas de cada solicitação. O cumprimento de metas não é um contador incrementado a cada doação, tampouco um campo persistido: é o resultado de um algoritmo de alocação FIFO calculado em memória a partir de um \textit{snapshot} do banco no momento da leitura.

\subsection{Recomendação de solicitações para o doador}

O caminho de leitura é implementado pelo \codigo{GetRecommendedRequestsUseCase}, acionado pelo \textit{endpoint} \codigo{GET /donation-requests/recommendations}. Dado um doador autenticado, o caso de uso primeiro verifica sua elegibilidade para doar na data corrente --- método \codigo{Donor.isEligibleToDonate(...)} (Listagem~\ref{lst:donor-eligible}) ---, considerando faixa etária e intervalo mínimo desde a última doação. Em seguida, busca no repositório as solicitações ativas compatíveis com o tipo sanguíneo do doador (Listagem~\ref{lst:bloodtype-candonateto}), dentro da distância máxima de recomendação configurada por ele (RF21) e ainda não expiradas, aproveitando o índice geoespacial descrito na Seção anterior.

As solicitações candidatas cuja meta já foi atingida no \textit{snapshot} de preenchimento são descartadas, e as demais são ordenadas por proximidade geográfica, urgência e proximidade da data-limite, de modo que o doador visualize primeiro as solicitações mais próximas, mais urgentes e mais próximas de expirar. O resultado retornado inclui, para cada solicitação recomendada, a distância estimada, a urgência, a meta de bolsas e o total já cumprido --- este último obtido do \codigo{DonationRequestFulfillmentService} no momento da leitura, sem campo persistido.

\subsection{Notificação de doadores potenciais}

Complementarmente à recomendação, o requisitante pode disparar notificações ativas por meio do \textit{endpoint} \codigo{POST /donation-requests/\{id\}/notify}, implementado pelo \codigo{NotifyPotentialDonorsUseCase}. Esse caso de uso valida que o solicitante autenticado é o dono da solicitação, que a solicitação ainda não expirou e que sua meta ainda não foi atingida; em seguida, delega ao \codigo{FindEligibleDonorsUseCase} (por meio do serviço de domínio \codigo{DonorMatchingService}, Listagem~\ref{lst:matching}) a identificação dos doadores elegíveis --- considerando compatibilidade sanguínea, elegibilidade temporal e proximidade geográfica --- para, então, notificar cada doador elegível com conta ativa através do serviço de notificação descrito na Seção~5.5.3. Esse fluxo materializa, na arquitetura da API, o comportamento observado na literatura (Capítulo~3) de notificar diretamente candidatos compatíveis em vez de aguardar passivamente novas doações.

\subsection{Snapshot do cumprimento de metas}
\label{sec:fulfillment-bags}

Não existe vínculo direto entre uma doação e uma solicitação específica: cada doação concluída em um hemocentro entra em um agrupamento (\textit{pool}) de doações daquela organização e é distribuída entre as solicitações ativas do mesmo hemocentro pelo \codigo{DonationRequestFulfillmentService} (Listagem~\ref{lst:fifo}). Os casos de uso de leitura invocam \codigo{fill} com a data de referência \codigo{asOfDate} --- em geral a data corrente. O serviço recarrega as solicitações ativas e ainda não expiradas dos hemocentros envolvidos (\codigo{findActiveRequestsByOrganizationIds}), de modo que a alocação considere o pool inteiro do centro, e não apenas o recorte da recomendação ou da listagem do requisitante. As doações concluídas são lidas na janela \texttt{[data da solicitação ativa mais antiga, asOfDate]}. A alocação ocorre em memória e isolada por hemocentro, com solicitações ordenadas da mais antiga para a mais nova e doações ordenadas pela data efetiva.

A política de distribuição combina a ordem \textit{First In, First Out} (FIFO) com um teto proporcional ao prazo restante de cada solicitação. O teto, calculado por \codigo{DonationRequest.proportionalGoalAt} (Listagem~\ref{lst:teto-proporcional}), libera a meta no mesmo ritmo em que a janela \codigo{[dateRequested, dateLimit]} é consumida: decorridos 3 de 10 dias, apenas $3/10$ da meta está disponível na data de referência. A distribuição ocorre, então, em duas passagens sobre o mesmo pool. Na primeira passagem, cada doação (com saldo inicial unitário contínuo de $1{,}0$ bolsa) é distribuída fracionariamente entre as solicitações compatíveis (\codigo{acceptsDonation}) até o respectivo teto proporcional; atingido o teto de um pedido, o saldo remanescente da doação avança na fila FIFO para as próximas solicitações compatíveis. Na segunda passagem, as frações de doações que nenhuma solicitação pôde absorver retornam na mesma ordem FIFO, agora alocadas até a meta cheia de cada solicitação.

O propósito do teto é redistribuir a escassez em favor de quem tem menos tempo. Sem ele, a solicitação mais antiga do hemocentro consome o pool integralmente, ainda que disponha de semanas de prazo, enquanto uma solicitação prestes a expirar permanece zerada. Como a segunda passagem devolve as sobras, a política não reduz o total de bolsas alocadas --- altera somente qual solicitação recebe cada bolsa quando o pool é insuficiente; havendo doações em quantidade suficiente para todas as metas, o resultado coincide com o da alocação FIFO simples.

Para ilustrar o funcionamento prático dessa estratégia sob uma ótica estritamente matemática contínua, a Figura~\ref{fig:exemplo-alocacao-teto} esquematiza um cenário hipotético em um mesmo hemocentro, avaliado na data de referência (\codigo{asOfDate}) de 24/07/2026.

\begin{figure}[H]
  \centering
  \caption{Cenário ilustrativo de alocação contínua restrita por teto proporcional}
  \includegraphics[width=0.98\textwidth]{figura10-exemplo-alocacao-teto}
  \fonte{Elaborado pelos próprios autores (2026).}
  \label{fig:exemplo-alocacao-teto}
\end{figure}

No cenário apresentado, existem duas solicitações ativas concorrendo pelo mesmo \textit{pool} de doações:
\begin{enumerate}[label=\alph*)]
  \item \textbf{Solicitação $R_1$ (mais antiga):} registrada em 14/07 com data-limite em 03/08 (janela total de 20 dias) e meta de 1,0 bolsa. Em 24/07, decorreram 10 dias da janela, correspondendo a uma fração decorrida de $10/20 = 0{,}5$ ($50\%$). Seu teto proporcional calculado é de $1{,}0 \times 0{,}5 = \mathbf{0{,}5000\text{ bolsa}}$;
  \item \textbf{Solicitação $R_2$ (mais recente e mais próxima do limite):} registrada em 21/07 com data-limite em 31/07 (janela total de 10 dias) e meta de 2,0 bolsas. Em 24/07, decorreram 3 dias da janela, correspondendo a uma fração decorrida de $3/10 = 0{,}3$ ($30\%$). Seu teto proporcional calculado é de $2{,}0 \times 0{,}3 = \mathbf{0{,}6000\text{ bolsa}}$.
\end{enumerate}

O hemocentro dispõe de um \textit{pool} com duas doações concluídas compatíveis: $D_1$ (concluída em 22/07, saldo de $1{,}0$ bolsa) e $D_2$ (concluída em 23/07, saldo de $1{,}0$ bolsa). 

Por se tratar de um \textit{snapshot} de alocação dinâmica calculado em memória para mensuração de desempenho puramente matemática --- sem persistência ou associação física estática prévia entre doações e receptores --- o modelo adota a soberania matemática contínua, deixando a distribuição a cargo da matemática sem arredondamentos artificiais:
\begin{itemize}
  \item \textbf{Metas e tetos proporcionais contínuos:} calculados em \codigo{BigDecimal}, permitindo que uma meta liberada temporalmente contenha frações de bolsa (por exemplo, $0{,}5$ ou $1{,}5$ bolsa);
  \item \textbf{Alocação dinâmica contínua:} o quantitativo alocado de uma doação a uma solicitação (\codigo{fulfilledBloodBags}) também pode conter frações de bolsa (por exemplo, $0{,}5$ ou $0{,}7$ bolsa), refletindo com exatidão a parcela da demanda atendida.
\end{itemize}

O fluxo de distribuição em duas passagens opera, portanto, da seguinte forma:
\begin{enumerate}
  \item \textbf{1ª Passagem (Teto Proporcional Contínuo):}
  \begin{itemize}
    \item \textbf{Distribuição de $D_1$ ($1{,}0$ bolsa):} a fila FIFO inicia pela solicitação mais antiga ($R_1$), cuja capacidade restante para atingir o teto é de $0{,}5 - 0{,}0 = 0{,}5$ bolsa. Aloca-se $0{,}5$ bolsa de $D_1$ para $R_1$, que atinge seu teto proporcional ($0{,}5000$). O saldo remanescente de $D_1$ ($1{,}0 - 0{,}5 = 0{,}5$ bolsa) avança na FIFO para $R_2$, que possui capacidade de $0{,}6 - 0{,}0 = 0{,}6$ bolsa. $R_2$ absorve integralmente essa fração restante de $0{,}5$ bolsa, esgotando o saldo de $D_1$ ($0{,}0$);
    \item \textbf{Distribuição de $D_2$ ($1{,}0$ bolsa):} $R_1$ já se encontra no limite do seu teto proporcional (capacidade zero), sendo represada na 1ª passagem. A busca avança para $R_2$, que possui capacidade restante de $0{,}6 - 0{,}5 = 0{,}1$ bolsa. Aloca-se $0{,}1$ bolsa de $D_2$ para $R_2$, que atinge seu teto proporcional ($0{,}6000$). O saldo não absorvido de $D_2$ ($1{,}0 - 0{,}1 = 0{,}9$ bolsa) não encontra mais capacidades disponíveis na 1ª passagem e é encaminhado à 2ª passagem.
  \end{itemize}
  \item \textbf{2ª Passagem (Meta Cheia / FIFO das Sobras):}
  \begin{itemize}
    \item Com os tetos temporais suspensos, aplicam-se as metas cheias ($1{,}0$ para $R_1$ e $2{,}0$ para $R_2$). A sobra de $0{,}9$ bolsa de $D_2$ retorna ao início da fila FIFO;
    \item $R_1$ necessita de $1{,}0 - 0{,}5 = 0{,}5$ bolsa para alcançar sua meta total. Aloca-se $0{,}5$ bolsa da sobra para $R_1$, que conclui integralmente sua meta ($1{,}0/1{,}0$ bolsa, \codigo{goalReached = true});
    \item O restante da sobra ($0{,}9 - 0{,}5 = 0{,}4$ bolsa) avança para $R_2$, elevando seu quantitativo final para $0{,}6 + 0{,}4 = 1{,}0$ bolsa ($50\%$ da meta cheia solicitada).
  \end{itemize}
\end{enumerate}

Ao término do processo, $R_1$ atinge $1{,}0/1{,}0$ bolsa ($100\%$ da meta) e $R_2$ atinge $1{,}0/2{,}0$ bolsas ($50\%$ da meta), com zero desperdício de doações. Nota-se o benefício dessa abordagem matemática: caso o hemocentro dispusesse de apenas uma doação no \textit{pool} ($D_1$ de $1{,}0$ bolsa), a doação seria partilhada exatamente como $0{,}5$ bolsa para $R_1$ e $0{,}5$ bolsa para $R_2$, impedindo que a solicitação mais antiga monopolizasse a única bolsa disponível e permitindo que ambas avançassem proporcionalmente em direção às suas metas.

O progresso (\codigo{fulfilledBloodBags}) não é persistido. Os casos de uso de leitura que precisam exibir o preenchimento --- listagem de solicitações do requisitante, recomendação ao doador e bloqueio de notificação quando a meta já foi atingida --- invocam o mesmo serviço de domínio e recebem o mapa calculado na hora. Concluir uma doação apenas persiste a doação; a leitura seguinte já vê o \textit{pool} atualizado. Como o cálculo não regrava solicitações, não há condição de corrida entre escritas de contador: cada leitura monta o próprio \textit{snapshot}. Assume-se, como contrapartida, que o progresso exibido de uma solicitação de prazo folgado pode decrescer entre duas leituras: à medida que os tetos das solicitações concorrentes crescem com o tempo, bolsas antes atribuídas a ela na segunda passagem migram para solicitações mais próximas do vencimento.
